\chapter{Requirements and Implementations}\label{ch:lists}

A project has an actual list of requirements, and a similar list for each of its
implementations. Both are lists of things that convert to CTL. An entry that does not
convert yet says so, the way \texttt{sorry} does in a Lean proof: the file still reads,
the entry is counted and warned about, and nothing pretends it holds.

This chapter follows one small shop from requirements to an implementation. Both files
are in \texttt{examples/orders}, and every verdict below is LeanFM's.

\section{What is required}

A requirements file says what is required in two ways. First, by a protocol: the
conversation the system must hold, message by message. This is PlaceOrder from
Figure~\ref{fig:placeorder}:

\begin{lstlisting}[language={},basicstyle=\ttfamily\footnotesize,breaklines=true]
/-- What the client orders. -/
structure CreateOrder where
  item : String
  quantity : Nat

/-- The shop's price for the order, in cents. -/
structure Quote where
  price : Nat

/-- The client pays; the card number is a secret, never logged. -/
structure Pay where
  amount : Nat
  card_secr : String

structure Charge where
  amount : Nat
  card_secr : String

structure Receipt where
  amount : Nat

/-- What is required: the conversation between a client, a shop and a bank, message by message. -/
protocol PlaceOrder (Client, Shop, Bank) {
  o: CreateOrder from Client to Shop route "POST /orders";
  choice {
    q: Quote from Shop to Client cause o status 201;
    p: Pay from Client to Shop cause q where p.amount = q.price;
    c: Charge from Shop to Bank cause p;
    choice {
      d: Charged from Bank to Shop cause c;
      Receipt from Shop to Client cause d status 200;
    } or {
      x: Declined from Bank to Shop cause c;
      PaymentDeclined from Shop to Client cause x status 402;
    }
  } or {
    SoldOut from Shop to Client cause o status 409;
  }
}
\end{lstlisting}


A client orders; the shop either quotes or says it is sold out; a quoted order is paid,
the bank charges the card, and the client gets a receipt or a refusal. The payment must
match the quote. Nothing else may happen, and every one of these paths must stay possible.

Second, by properties the protocol's shape does not show at a glance. The file ends in a
map that names every requirement and points at what it means:

\begin{lstlisting}[language={},basicstyle=\ttfamily\footnotesize,breaklines=true]
-- from examples/orders/Requirements.leanfm
requirements {
  "Orders are placed as PlaceOrder says" => PlaceOrder;
  "Every order is answered" => PlaceOrder ⊨ (leadsTo (on CreateOrder) (or (on Receipt) (on PaymentDeclined) (on SoldOut)));
  "Every payment is charged to the bank" => PlaceOrder ⊨ (leadsTo (on Pay) (on Charge));
  "No receipt without a charge" => PlaceOrder ⊨ (precedes (on Charged) (on Receipt));
  "A sold-out order is never charged" => PlaceOrder ⊨ (always (implies (on SoldOut) (never (on Charge))));
  "Refunds within 30 days" => sorry "PlaceOrder has no Refund message yet";
}
\end{lstlisting}


Each entry converts to CTL and is decided. Here is what each one requires, as the CTL
LeanFM writes for it and as LeanFM reads that formula back, with the verdict:

\begin{lstlisting}[language={},basicstyle=\ttfamily\footnotesize,breaklines=true]
"Orders are placed as PlaceOrder says"
    the conversation is exactly PlaceOrder's: 11 faithful CTL rules, one per typestate and one for the start
    -> ok

"Every order is answered"
    AG (CreateOrder → AF (Receipt ∨ PaymentDeclined ∨ SoldOut))
    on every path, always (if CreateOrder was just delivered then (on every path, eventually (Receipt was just delivered or PaymentDeclined was just delivered or SoldOut was just delivered)))
    -> proved

"Every payment is charged to the bank"
    AG (Pay → AF Charge)
    on every path, always (if Pay was just delivered then (on every path, eventually Charge was just delivered))
    -> proved

"No receipt without a charge"
    ¬E[¬Charged U Receipt]
    it is not the case that (on some path, (it is not the case that Charged was just delivered) holds until Receipt was just delivered (and Receipt was just delivered does happen))
    -> proved

"A sold-out order is never charged"
    AG (SoldOut → AG ¬Charge)
    on every path, always (if SoldOut was just delivered then (on every path, always (it is not the case that Charge was just delivered)))
    -> proved

"Refunds within 30 days"
    sorry: PlaceOrder has no Refund message yet
    -> sorry
\end{lstlisting}


A pointer is one of three kinds:

\begin{center}
\begin{tabular}{@{}lp{.62\textwidth}@{}}
\toprule
pointer & converts to \\
\midrule
\texttt{P} & the faithful rules of protocol $P$ (Chapter~\ref{ch:faithful}): exactly $P$'s conversation \\
\texttt{P} $\models$ \texttt{($\varphi$)} & the formula $\varphi$, decided over the state graph of $P$ \\
\texttt{sorry} & nothing yet \\
\bottomrule
\end{tabular}
\end{center}

A pointer to a protocol the file does not define is an error, not a sorry, as an
unknown name is in Lean. A name appears once, and the map is the file's last item, so a
reader always knows where the requirements are. The map changes no protocol.

\section{Sorry}

\texttt{sorry} is anything that does not convert to CTL yet. ``Refunds within 30 days''
cannot: PlaceOrder has no refund message, so there is nothing for a formula to talk
about. Writing the requirement down anyway, as a sorry, keeps it on the list.

A sorry is neither true nor false. Converting it to \emph{true} would count an unwritten
requirement as met; converting it to \emph{false} would claim it is refuted. Like Lean's
\texttt{sorry}, it is admitted and tracked: the check passes with a warning for each one,
and in strict mode each one is an error. Removing a sorry means adding to the protocol, here
a Refund message and a deadline, and replacing the sorry with a formula over it. A pointer
to a kind with no transform yet, such as an SLA block, whose claims are probabilistic, is
a sorry too.

\section{An implementation}

Some things belong in a specification but not in its requirements: behaviour no client
can observe, and choices that are not behaviour at all, such as the language or the
database. They go in an implementation, which names the requirements it implements. A
project may have any number of implementations against the same requirements.

The Go shop puts an API in front of a database. The database is invisible to the client,
so its protocol, ShopImpl, lives in the implementation:

\begin{lstlisting}[language={},basicstyle=\ttfamily\footnotesize,breaklines=true]
protocol ShopImpl (Client, Api, Db, Bank) {
  o: CreateOrder from Client to Api route "POST /orders";
  r: Reserve from Api to Db cause o;
  choice {
    v: Reserved from Db to Api cause r;
    q: Quote from Api to Client cause o status 201;
    p: Pay from Client to Api cause q;
    c: Charge from Api to Bank cause p;
    choice {
      d: Charged from Bank to Api cause c;
      w: RecordPayment from Api to Db cause d;
      k: Recorded from Db to Api cause w;
      Receipt from Api to Client cause d status 200;
    } or {
      x: Declined from Bank to Api cause c;
      ReleaseStock from Api to Db cause x;
      PaymentDeclined from Api to Client cause x status 402;
    }
  } or {
    n: OutOfStock from Db to Api cause r;
    SoldOut from Api to Client cause o status 409;
  }
}
\end{lstlisting}


\begin{figure}[ht]\centering
\AccessibleFigure{Interaction diagram of ShopImpl. Message from Client to Api: CreateOrder. Message from Api to Db: Reserve. Then either: message from Db to Api: Reserved; message from Api to Client: Quote; message from Client to Api: Pay; message from Api to Bank: Charge; then either message from Bank to Api: Charged; message from Api to Db: RecordPayment; message from Db to Api: Recorded and message from Api to Client: Receipt or message from Bank to Api: Declined; message from Api to Db: ReleaseStock and message from Api to Client: PaymentDeclined. Or: message from Db to Api: OutOfStock; message from Api to Client: SoldOut.}{\ThemedDiagram{width=.85\textwidth}{req-shopimpl}}
\caption{ShopImpl: the shop as built, with the database the client never sees.}
\label{fig:shopimpl}
\end{figure}

LeanFM's projection check shaped this protocol: a first version told the database about
a payment only when the bank accepted it, and LeanFM refused it because the database could
not tell which way the bank had decided. ReleaseStock is the fix.

The implementation's block holds three kinds of entry:

\begin{lstlisting}[language={},basicstyle=\ttfamily\footnotesize,breaklines=true]
-- from examples/orders/Implementation.leanfm
implementation Go of examples.orders.Requirements {
  refines PlaceOrder by ShopImpl;
  "Orders are placed as PlaceOrder says" => conform "examples/orders/logs/order-1.json";
  "Every order is answered" => test "shop/order_test.go: TestEveryOrderAnswered";
  "Every payment is charged to the bank" => conform "examples/orders/logs/order-1.json";
  "No receipt without a charge" => sorry;
  "Stock is reserved before a quote" => ShopImpl ⊨ (precedes (after Reserved) (after Quote));
  "A payment is recorded before the receipt" => ShopImpl ⊨ (precedes (after Recorded) (after Receipt));
  hint "Implement the code as Go, with net/http and no web framework";
  hint "Keep orders in PostgreSQL; the Db role is one transaction per message";
}
\end{lstlisting}


\begin{description}
\item[Evidence] for a requirement: \texttt{conform "log"} checks a run's log against the
requirement's protocol (Chapter~\ref{ch:where}); \texttt{test "name"} names a test the
project's checks run; \texttt{sorry} has none yet. A name the requirements do not have is
an error.
\item[Its own CTL], under new names: here, that stock is reserved before a quote and a
payment recorded before the receipt. Both are about the database, which is why they are
not requirements. \texttt{refines PlaceOrder by ShopImpl} decides every formula requirement
over PlaceOrder again over ShopImpl, so the implementation may add messages but may not
break what was required.
\item[Hints] for whoever writes the code, often an LLM: the language, the database. A
hint claims no behaviour, so it converts to \emph{true}, honestly, and never warns.
\end{description}

What each entry means and how it stands:

\begin{lstlisting}[language={},basicstyle=\ttfamily\footnotesize,breaklines=true]
refines PlaceOrder by ShopImpl
    every formula requirement over PlaceOrder decided again over ShopImpl: ShopImpl keeps all 4 formula requirements over PlaceOrder
    -> refines

"Orders are placed as PlaceOrder says"
    evidence: the log examples/orders/logs/order-1.json conforms to the requirement's protocol
    -> met

"Every order is answered"
    evidence: the test shop/order_test.go: TestEveryOrderAnswered
    -> met

"Every payment is charged to the bank"
    evidence: the log examples/orders/logs/order-1.json conforms to the requirement's protocol
    -> met

"No receipt without a charge"
    sorry
    -> specified

"Stock is reserved before a quote"
    precedes(Reserved, Quote)
    on every path, Quote was just delivered never holds unless Reserved was just delivered has held at or before that point (Quote was just delivered need not hold at all)
    -> proved

"A payment is recorded before the receipt"
    precedes(Recorded, Receipt)
    on every path, Receipt was just delivered never holds unless Recorded was just delivered has held at or before that point (Receipt was just delivered need not hold at all)
    -> proved

hint 1: Implement the code as Go, with net/http and no web framework
    claims no behaviour, converts to true
    -> true

hint 2: Keep orders in PostgreSQL; the Db role is one transaction per message
    claims no behaviour, converts to true
    -> true
\end{lstlisting}


\section{Status is computed}

Status follows from the two verdicts, so no one has to declare it:

\begin{center}
\begin{tabular}{@{}llll@{}}
\toprule
requirement & evidence & status & the check \\
\midrule
converts and holds & passes & met & ok \\
converts and holds & sorry or none & specified & warning \\
sorry & passes & tested & warning \\
sorry & sorry or none & sorry & warning \\
fails & any & failing & error \\
any & fails & failing & error \\
\bottomrule
\end{tabular}
\end{center}

A requirement an implementation does not name counts as sorry for it. The shop:

\begin{lstlisting}[language={},basicstyle=\ttfamily\footnotesize,breaklines=true]
examples/orders: 15 entries; 0 failing; 3 use sorry
requirements: 5 convert to CTL, 5 met
implementation Go: 2 proved, 0 tested, 4 sorry, 0 failing
\end{lstlisting}


Evidence is only as strong as its kind. A conforming log carries every $\AG$ property
proved of the protocol to that run, since the run is one of the protocol's runs; an
$\AF$ or $\EF$ property also needs the guards of every choice to cover every value
(Chapter~\ref{ch:where}). A test is evidence but not CTL, so a sorry requirement with a
passing test is ``tested'', not ``met''.

\section{LeanFM's own lists}

LeanFM is specified the same way. Its requirements map ends
\texttt{LeanFM/LLMGenerated/LeanFM.leanfm}, over its own MCP protocol; its implementation
is \texttt{LeanFMImplementation.leanfm}; and the clauses of \texttt{SelfHosting.lean},
which are still prose, join its list as sorry entries:

\begin{lstlisting}[language={},basicstyle=\ttfamily\footnotesize,breaklines=true]
LeanFM: 147 entries; 0 failing; 131 use sorry
requirements (LeanFM/LLMGenerated/LeanFM.leanfm): 7 convert to CTL, 7 met; 129 use sorry (the clauses of SelfHosting.lean)
implementation Lean (LeanFM/LLMGenerated/LeanFMImplementation.leanfm): 5 met, 2 specified, 0 tested, 0 sorry, 0 failing
\end{lstlisting}


\section{The gate and the viewers}

\texttt{make requirements-check} converts and decides every entry. A failing entry fails
it; each sorry is a warning, and with \texttt{STRICT=1} an error. It checks LeanFM's
specification directory and, with \texttt{--files}, any other files, such as this
chapter's example.

The session view lists both lists under \emph{spec}, \emph{files \& sections}. Each entry
opens in the viewer of its kind: a protocol in the protocol explorer, a formula in the CTL
view with its witness or counterexample path, and a sorry, a hint or a piece of evidence as
text. A new kind of entry brings its transform and its viewer; until it has a transform it
is a sorry.
